\documentclass[12pt]{article}

\usepackage[a4paper,margin=1in]{geometry}
\usepackage{amsmath,amssymb,mathtools}
\usepackage{graphicx}
\usepackage{booktabs}
\usepackage{enumitem}
\usepackage{setspace}
\usepackage{hyperref}
\usepackage{caption}

\hypersetup{
    colorlinks=true,
    linkcolor=blue,
    urlcolor=blue,
    citecolor=blue
}

\setstretch{1.15}

\title{A Mathematical Report on Planetary Motion Simulations}
\author{Lee Cheuk Man Gerard}
\date{}

\begin{document}

\maketitle

\begin{abstract}
This report presents a mathematical analysis of two planetary motion visualizations:
\textit{planetary\_motion\_kepler\_2d.mp4} and \textit{planetary\_motion\_nbody.mp4}.
The first output illustrates the classical Keplerian description of planetary orbits,
emphasizing elliptical trajectories and non-uniform orbital speed.
The second output implements a Newtonian $N$-body gravitational model in which the
motion of each body arises from mutual interactions with all other bodies.
Together, the two simulations highlight the distinction between a kinematic orbital
approximation and a dynamical many-body system governed by differential equations.
\end{abstract}

\section{Kepler's Laws of Planetary Motion}

The output \textit{planetary\_motion\_kepler\_2d.mp4} provides a geometric visualization
of planetary orbits in the heliocentric plane. In this representation, each planet
moves along an ellipse with the Sun located at one focus, in accordance with
Kepler's first law. The animation also demonstrates the non-uniform orbital speed
predicted by Kepler's second law and the period--semimajor axis relationship of
Kepler's third law.

\subsection{Visual Output}

Figure~\ref{fig:kepler-output} shows the representative frame extracted from the
\textit{planetary\_motion\_kepler\_2d.mp4} output. The image file
\texttt{kepler\_output.png} should be replaced with the corresponding frame exported
from the video.

\begin{figure}[h!]
    \centering
    \includegraphics[width=0.92\textwidth]{kepler_output.png}
    \caption{Representative visual output from \textit{planetary\_motion\_kepler\_2d.mp4}.}
    \label{fig:kepler-output}
\end{figure}

\subsection{Orbital Geometry}

Let a planet have semimajor axis \(a\) and eccentricity \(e\). In its orbital plane,
the trajectory may be parameterized by the eccentric anomaly \(E\) as
\begin{equation}
x'(E) = a(\cos E - e), 
\qquad
y'(E) = a\sqrt{1-e^2}\sin E.
\end{equation}
This parameterization describes an ellipse with the central body located at a focus.

The motion in time is governed by Kepler's equation,
\begin{equation}
M = E - e\sin E,
\end{equation}
where \(M\) denotes the mean anomaly. Since
\begin{equation}
M(t) = M_0 + nt,
\end{equation}
with mean motion \(n\), the orbital position evolves nonlinearly in time even
though the mean anomaly increases uniformly.

\subsection{Kepler's Second Law}

Kepler's second law states that the radius vector from the Sun to the planet sweeps
out equal areas in equal times. In differential form, the areal velocity is constant:
\begin{equation}
\frac{dA}{dt} = \frac{1}{2}r^2\frac{d\theta}{dt} = \text{constant}.
\end{equation}
Equivalently, over an infinitesimal displacement \(d\mathbf{r}\), the swept area is
\begin{equation}
dA = \frac{1}{2}\left|\mathbf{r} \times d\mathbf{r}\right|.
\end{equation}
This property is visible in the animation through the varying orbital speed:
the planets move faster near perihelion and slower near aphelion.

\subsection{Kepler's Third Law}

Kepler's third law gives the relation between the orbital period \(T\) and the
semimajor axis \(a\):
\begin{equation}
T^2 \propto a^3.
\end{equation}
In the commonly used normalized astronomical system, this relation is written as
\begin{equation}
T^2 = a^3.
\end{equation}
Accordingly, planets with larger orbital radii exhibit longer periods, which is
consistent with the visual differences observed in the simulation.

\subsection{Interpretation of the Output}

The output \textit{planetary\_motion\_kepler\_2d.mp4} is best interpreted as a
kinematic model of planetary motion rather than a force-based numerical simulation.
Its principal mathematical features are summarized as follows:

\begin{enumerate}[label=\arabic*.]
    \item The orbit of each planet is elliptical.
    \item The Sun remains fixed at one focus of the ellipse.
    \item The orbital speed varies along the trajectory.
    \item Equal areas are swept out in equal times.
    \item The orbital period satisfies \(T^2 \propto a^3\).
\end{enumerate}

Thus, the animation serves as an effective visualization of the classical laws
describing planetary motion.

\section{Newtonian Gravity Simulation}

The output \textit{planetary\_motion\_nbody.mp4} presents a Newtonian gravitational
simulation in which the motion of each body is determined by the combined influence
of all other bodies in the system. Unlike the Keplerian output, this animation does
not prescribe orbital paths in advance; instead, trajectories emerge from the direct
solution of the equations of motion.

\subsection{Visual Output}

Figure~\ref{fig:nbody-output} shows a representative frame extracted from the
\textit{planetary\_motion\_nbody.mp4} output. The image file
\texttt{nbody\_output.png} should be replaced with the corresponding frame exported
from the video.

\begin{figure}[h!]
    \centering
    \includegraphics[width=0.92\textwidth]{nbody_output.png}
    \caption{Representative visual output from \textit{planetary\_motion\_nbody.mp4}.}
    \label{fig:nbody-output}
\end{figure}

\subsection{Equations of Motion}

Let \(\mathbf{r}_i(t)\) be the position vector of body \(i\) with mass \(m_i\).
Then the Newtonian \(N\)-body equations are
\begin{equation}
\ddot{\mathbf{r}}_i
=
G \sum_{\substack{j=1 \\ j\neq i}}^{N}
m_j \frac{\mathbf{r}_j - \mathbf{r}_i}{\left|\mathbf{r}_j - \mathbf{r}_i\right|^3},
\qquad i=1,\dots,N.
\end{equation}
This coupled system of second-order differential equations describes the gravitational
acceleration of each body due to all the others.

The center of mass is defined by
\begin{equation}
\mathbf{R}_{\mathrm{cm}}
=
\frac{1}{M}\sum_{i=1}^{N} m_i \mathbf{r}_i,
\qquad
M = \sum_{i=1}^{N} m_i.
\end{equation}
If the total momentum is initially zero, the center of mass remains fixed.
This explains the barycentric motion of the Sun in the animation.

\subsection{Numerical Integration}

For numerical propagation, a symplectic scheme such as velocity Verlet is commonly
used. Given positions \(\mathbf{r}_i^n\), velocities \(\mathbf{v}_i^n\), and
accelerations \(\mathbf{a}_i^n\) at time step \(n\), the update is
\begin{equation}
\mathbf{r}_i^{n+1}
=
\mathbf{r}_i^n + \Delta t\,\mathbf{v}_i^n + \frac{1}{2}(\Delta t)^2 \mathbf{a}_i^n,
\end{equation}
followed by recomputation of \(\mathbf{a}_i^{n+1}\), and then
\begin{equation}
\mathbf{v}_i^{n+1}
=
\mathbf{v}_i^n + \frac{\Delta t}{2}\left(\mathbf{a}_i^n + \mathbf{a}_i^{n+1}\right).
\end{equation}
This method is particularly suitable for orbital simulations because it preserves
the qualitative structure of the dynamics over long time intervals.

\subsection{Interpretation of the Output}

The output \textit{planetary\_motion\_nbody.mp4} is a direct numerical realization
of Newtonian gravitational mechanics. Its main features may be summarized as follows:

\begin{enumerate}[label=\arabic*.]
    \item Each body experiences the combined gravitational force of all other bodies.
    \item The Sun is not stationary, but moves slightly about the barycenter.
    \item The resulting trajectories are approximately elliptical but not exactly closed.
    \item Small perturbations arise from mutual planetary interactions.
\end{enumerate}

Consequently, this animation provides a more physically complete description of
planetary dynamics than the Keplerian approximation, while also demonstrating the
computational complexity of the many-body problem.

\begin{thebibliography}{9}

\bibitem{kepler1609}
J.~Kepler.
\newblock \textit{Astronomia Nova}.
\newblock Prague, 1609.

\bibitem{newton1687}
I.~Newton.
\newblock \textit{Philosophi\ae{} Naturalis Principia Mathematica}.
\newblock London, 1687.

\bibitem{goldstein}
H.~Goldstein, C.~Poole, and J.~Safko.
\newblock \textit{Classical Mechanics}, 3rd ed.
\newblock Pearson, 2002.

\bibitem{bate}
R.~R. Bate, D.~D. Mueller, and J.~E. White.
\newblock \textit{Fundamentals of Astrodynamics}.
\newblock Dover Publications, 1971.

\end{thebibliography}

\end{document}